\documentclass{beamer}

\usepackage[utf8]{inputenc}
\usepackage[T1]{fontenc}
\usepackage[francais]{babel}
\usepackage{lmodern}
\usepackage{color}
\usepackage{fix-cm}
\usepackage{textpos}
\usepackage{eurosym}
\usepackage{multirow}
\usepackage{perpage}

% Redéfinis les marges des tableaux
\let\oldtabular=\tabular
\def\tabular{\small\oldtabular}
\renewcommand{\arraystretch}{1.5}

\usetheme{Warsaw}
\usecolortheme{orchid}

% Permet de réinitialiser les footnote à chaque frame. Nécéssite 2 compilations.
\MakePerPage{footnote}

\setlength{\TPHorizModule}{0.01\textwidth}
\setlength{\TPVertModule}{0.01\textheight}

\setbeamertemplate{navigation symbols}{}

\title[Réseaux sociaux et communication]{
    La communication via les réseaux sociaux}
\author[HOUDAYER \and RUHIER \and SANTIAGO \and TSCHIRRET]{
	Benoit HOUDAYER\\
    \and
	Anthony RUHIER\\
    \and
    Cédric SANTIAGO\\
    \and
    Maxime TSCHIRRET
}
\institute[TI81 - UTBM]{
    Communication (TI81) - Université Technologique de Belfort Montbéliard}
\date{Avril 2015}

\begin{document}

% Permet de masquer le comptage de slides
\bgroup
\makeatletter
\setbeamertemplate{headline}{}
\setbeamertemplate{footline}
{
  \leavevmode%
  \hbox{%
  \begin{beamercolorbox}[wd=.5\paperwidth,ht=2.25ex,dp=1ex,center]{title in head/foot}%
    \usebeamerfont{title in head/foot}\insertshorttitle
  \end{beamercolorbox}%
  \begin{beamercolorbox}[wd=.5\paperwidth,ht=2.25ex,dp=1ex,center]{date in head/foot}%
    \usebeamerfont{date in head/foot}\insertshortdate{}
%    \insertframenumber{} / \inserttotalframenumber\hspace*{2ex}
  \end{beamercolorbox}}%
}
\makeatother

\begin{frame}[noframenumbering]
 	\frametitle{}
	\titlepage
\end{frame}
\egroup

\begin{frame}[noframenumbering,plain]
\begin{center}
    \begin{textblock}{90}(-35,-30)
        \includegraphics[width=25mm]{images/facebook_logo.eps}
    \end{textblock}

    \begin{textblock}{90}(45,-30)
        \includegraphics[width=20mm]{images/reddit_logo.eps}
    \end{textblock}

    \vspace{4em}
    \bf{\Large{Qu'est ce qu'un réseau social ?}}

    \begin{textblock}{90}(45,21)
        \includegraphics[width=25mm]{images/twitter_logo.eps}
    \end{textblock}

    \begin{textblock}{90}(-35,18)
        \includegraphics[width=25mm]{images/github_logo.eps}
    \end{textblock}
\end{center}
\end{frame}

\logo{\includegraphics[height=1cm]{images/logo-utbm.eps}\hspace*{2em}}


% Ajout du compteur de slides
\expandafter\def\expandafter\insertshorttitle\expandafter{%
      \insertshorttitle\hfill%
      \insertframenumber\,/\,\inserttotalframenumber}

\begin{frame}
	\begin{small}
	\frametitle{Sommaire}
	\tableofcontents
	\end{small}
\end{frame}

\AtBeginSection[]
{
	\begin{small}
	\begin{frame}
		\frametitle{Sommaire}
		\tableofcontents[currentsection]
	\end{frame}
	\end{small}
}


%%%% Includes des sections :
%%%%%%%%%%%%%%%%%%%%%%%%%%%%%%
%
    \section{Réseaux personnels}

\subsection{Présentation}
\begin{frame}
\frametitle{Présentation}
\begin{itemize}
    \itemsep1.5em
    \item Exemples :
        \begin{itemize}
            \itemsep0.5em
            \item Facebook
            \item Twitter
            \item Instagram
        \end{itemize}
    \item Quelques chiffres :
        \begin{itemize}
            \itemsep0.5em
            \item 75\% de la population mondiale sur les réseaux sociaux % vérifier chiffres
            \item 68\% des personnes en France
            \item 1h30 d'utilisation/jour
        \end{itemize}
\end{itemize}
\end{frame}

\subsection{Objectifs}
\begin{frame}
\frametitle{Objectifs}
\begin{itemize}
    \itemsep1.5em
    \item Pour la personne physique
        \begin{itemize}
            \itemsep0.5em
            \item Se divertir
            \item Communiquer
            \item Échanger, partager
        \end{itemize}
    \item Pour la personnes morale
        \begin{itemize}
            \itemsep0.5em
            \item Développement d'images de marque
            \item Communication avec les consommateurs
            \item Surveiller son e-reputation et celle de ses concurrents
            \item Une meilleure visibilité pour un moindre coût
        \end{itemize}
\end{itemize}
\end{frame}

% Ajouter 3ème section

    \section{Réseaux de contenu}

\subsection{Présentation}
\begin{frame}
\frametitle{Présentation}
\begin{itemize}
    \itemsep1em
    \item Exemples : Reddit, Slashdot, Digg
    \item Utilisateurs anonymes
    \item Liens inter-utilisateurs au second plan
    \item Centré sur le contenu
    \item Regroupement par thème
\end{itemize}
\end{frame}

\subsection{Une communauté plus qu'une personne}
\begin{frame}
\frametitle{Une communauté plus qu'une personne}
\begin{itemize}
    \itemsep1em
    \item Contenu sélectionné et hiérarchisé par les communautés
    \item Émergence de communautés
    \item Expression d'opinions
    \item Adjacence des communautés
\end{itemize}
\end{frame}

    \section{Réalisations/Projets}

\subsection{Mise en avant des réalisations}
\begin{frame}
\frametitle{Mise en avant des réalisations}
\begin{itemize}
    \itemsep1.5em
    \item Très peu d'informations sur l'équipe ou la personne
    \item Le jugement, la notation et les abonnements sur font sur les travaux
    \item Pousse au partage
    \item Se retrouver dans la masse de contenu
\end{itemize}
\end{frame}

\subsection{Une façon différente de travailler}
\begin{frame}
\frametitle{Une façon différente de travailler}
\begin{itemize}
    \itemsep2em
    \item La communauté peut suivre le travail en cours
    \item Influe sur le projet
    \item Le créateur peut mieux prévoir le succès de ses créations
\end{itemize}
\end{frame}

\subsection{Thématiques larges}
\begin{frame}
\frametitle{Thématiques larges}
\begin{itemize}
    \itemsep1.5em
    \item Majoritairement utilisé dans les métiers de l'informatique. Exemples:
        \begin{itemize}
            \itemsep0.3em
            \item Développement logiciel : Github
            \item Modélisation 3D
            \item Graphisme : Deviantart
        \end{itemize}
    \item Sites laissant l'utilisateur fournir le contenu et le critiquer
    \item Cuisine: proposition de recettes, commentaires et notes
\end{itemize}
\end{frame}


\end{document}
